% ============================================================
%  Chapter 6 -- Planning as Model Checking
% ============================================================

\chapter{Planning as Model Checking}
\label{ch:planning-model-checking}
\chaptermindmap{06}

\aside{Planning can be read as a
controller-synthesis problem: find a
strategy that drives the system to a
goal despite uncertainty.}

Part~I has built a toolbox: finite automata to recognise behaviours,
omega-automata for infinite traces, congruences to compare languages,
and games with attractors to synthesise controllers.  This closing
chapter changes lens.  We step outside verification for a moment and
look at \emph{automated planning}, the classical AI problem of choosing
actions that drive a world from an initial situation to a goal.  The
surprise, and the reason this chapter belongs here, is that planning
\emph{is} the same problem in disguise.  A planning domain is a finite
transition system, a goal is a target set of states, and a plan under
uncertainty is exactly a strategy in a game against nature.  The
preimage operator that defined the system attractor in \cref{ch:synthesis-games}
will reappear verbatim as the one-step engine of backward planning.

A planning problem describes a world, a set of actions, and a goal.
The task is to construct a plan that reaches the goal.  When actions
are non-deterministic or the environment can interfere, a plan is no
longer just a finite sequence of actions; it is a strategy.  We will
classify the strength of such strategies --- weak, strong, and
strong-cyclic --- and show that each corresponds to a familiar
verification flavour: existential reachability, universal reachability,
and reachability under fairness.  This recasting of planning as model
checking is what gives the chapter its name and what ties Part~I
together before we turn, in Part~II, to the temporal logics and
symbolic tooling that make these computations scale.  Model checking
is not the only view: as the professor pointed out, modern planning
also translates into satisfiability checking (SAT), controller synthesis
(building the state-machine that implements the plan), and execution
monitoring (detecting when the environment deviates from assumptions
and dynamically re-planning).

\begin{chapteropening}
\textbf{Where we are.}  The previous chapters built automata, language
operations, complementation, and games.  The key algorithmic pattern was
already visible in reachability games: compute the states from which a player
can force progress.

\textbf{What this chapter adds.}  Planning recasts that pattern as an AI
problem.  Actions change a finite world, uncertainty plays the role of an
adversary, and a plan is a strategy for reaching a goal.

\textbf{Where it leads.}  The same preimage and fixpoint ideas will return in
Part~II as symbolic model checking, CTL fixed points, and tool-supported
verification.

\textbf{Chapter map.}
\begin{itemize}[nosep]
  \item \Cref{sec:action-vs-timeline-planning} compares action-based and
        timeline-based modelling.
  \item \Cref{sec:planning-domains} formalises fluents, states, and actions.
  \item \Cref{sec:plans-execution-structures} explains plans as strategies.
  \item \Cref{sec:preimage-operators,sec:planning-fixpoints} develop the
        backward operators and algorithms.
  \item \Cref{sec:symbolic-planning,sec:action-description-languages} show how
        compact symbolic descriptions replace explicit graphs.
  \item \Cref{sec:planning-towards-mc} connects planning back to model
        checking and closes Part~I.
\end{itemize}
\end{chapteropening}

% ============================================================
\section{Action-Based and Timeline-Based Planning}
\label{sec:action-vs-timeline-planning}
% ============================================================

Before formalising anything, we have to choose how to describe the
world that is being controlled.  Two families of models have emerged
from decades of practice, and they emphasise different things.  One
looks at the world through the actions an agent performs; the other
looks at it through the values that state variables take over time.
The choice is not cosmetic: it shapes which kinds of constraint feel
natural to write and which algorithms fit.

\begin{definition}[Action-based planning]
\label{def:action-based-planning}
\textbf{Action-based planning} models change as a sequence of discrete
actions.  A plan is built from operations such as \texttt{load},
\texttt{move}, or \texttt{cross}, each with preconditions and effects.
\end{definition}

This is the graph view: states are nodes, actions are labelled
transitions, and a plan is a path or policy leading to a goal.  It is
the view that connects directly to the automata and games of the
previous chapters.

\begin{definition}[Timeline-based planning]
\label{def:timeline-based-planning}
\textbf{Timeline-based planning} models the evolution of state variables
over time intervals.  Instead of focusing on individual actions, it
records when each component has a given value and which temporal
constraints synchronize different components.
\end{definition}

Timeline planning is natural for scheduling and control domains.  For a
spacecraft, for instance, an instrument may be \emph{off}, then
\emph{warming up}, then \emph{collecting data}; another component, such
as a battery or data buffer, evolves on its own timeline but must remain
compatible with the instrument timeline.  This style was used for
real NASA missions and is closer in spirit to control theory than to
graph search.

\begin{figure}[ht]
\centering
\begin{tikzpicture}[
  slot/.style={draw, minimum height=7mm, align=center, font=\small},
]
  \node[anchor=east] at (0,0) {Instrument};
  \node[slot, fill=gray!10, minimum width=21mm] at (1.2,0) {off};
  \node[slot, fill=orange!12, minimum width=23mm] at (3.4,0) {warm up};
  \node[slot, fill=teal!12, minimum width=32mm] at (6.1,0) {collect};
  \node[anchor=east] at (0,-0.9) {Buffer};
  \node[slot, fill=gray!10, minimum width=39mm] at (2.1,-0.9) {empty};
  \node[slot, fill=violet!10, minimum width=37mm] at (5.9,-0.9) {filling};
  \draw[-Stealth, thick] (0,-1.6) -- (7.8,-1.6) node[right] {time};
\end{tikzpicture}
\caption{Timeline-based planning follows component values over intervals and synchronizes them by constraints.}
\label{fig:timeline-planning}
\end{figure}

The rest of this chapter focuses on the action-based view because it
connects directly to transition systems and model checking.  The
timeline view is important in applications, but it requires richer
temporal constraints.


% ============================================================
\section{Planning Domains}
\label{sec:planning-domains}
% ============================================================

We now commit to the action-based view and build the formal model.
Just as a finite automaton is a finite set of states with a labelled
transition relation, an action-based planning domain is a finite set
of states equipped with actions that take a state to one or more
successors.  The two ingredients we add beyond pure automata are an
\emph{initial} set (where the world might start) and a \emph{goal} set
(where we want it to end up).

\begin{definition}[Fluent and state]
\label{def:fluent-state}
A \textbf{fluent} is a Boolean variable describing a property of the
world.  A \textbf{state} is a valuation of all fluents.
\end{definition}

If the fluents are $\mathit{atA}$, $\mathit{loaded}$, and
$\mathit{green}$, then a state records whether each of these facts is
currently true.  With $n$ Boolean fluents, the explicit state space may
contain up to $2^n$ states --- our first hint that symbolic methods
will eventually be needed.

\begin{definition}[Inertial and non-inertial fluents]
\label{def:inertial-fluents}
An \textbf{inertial fluent} keeps its value unless some action or event
changes it.  A \textbf{non-inertial fluent} is recomputed or controlled
by the environment at each step.
\end{definition}

The split between inertial and non-inertial fluents is the planning
analogue of the system/environment split in \cref{ch:synthesis-games}.
The planner controls inertial fluents through its actions; non-inertial
fluents are decided by nature.  For the parcel domain below, the
robot's location is inertial, while the colour of the traffic light is
non-inertial from the robot's point of view.

\begin{definition}[Action-based planning domain]
\label{def:planning-domain}
An action-based planning domain consists of:
\begin{itemize}[nosep]
  \item a finite set of states $S$;
  \item a set of initial states $I\subseteq S$;
  \item a set of goal states $G\subseteq S$;
  \item a finite set of actions $A$;
  \item a transition relation $T\subseteq S\times A\times S$.
\end{itemize}
An action $a$ is \textbf{enabled} at $s$ when there exists at least one
$s'$ with $(s,a,s')\in T$.  The set of possible successors of $(s,a)$
is $\mathrm{Post}(s,a)=\{s'\in S:(s,a,s')\in T\}$.
\end{definition}

The transition relation may be non-deterministic: the same action in
the same state can have several possible successors.  This is where
planning begins to look like a game against the environment.  The
domain by itself fixes only the rules; the planner must still pick a
move at each turn, and nature resolves the residual non-determinism.

\begin{formaldetails}[Modelling dimensions]
Planning domains are often classified along three independent axes.
\begin{itemize}[nosep]
  \item \textbf{Deterministic vs. non-deterministic effects}: an action
        has one successor, or several possible successors.
  \item \textbf{Complete vs. partial observability}: the planner sees
        the exact current state, or only an observation compatible with
        several states.
  \item \textbf{Reachability vs. extended goals}: the goal is simply to
        reach $G$, or to satisfy temporal requirements along the way.
        \aside{In the lecture, the professor illustrated this with a driving
        analogy: the destination is the goal state, but during the trip
        you must also respect a velocity threshold --- a constraint that
        must hold \emph{along the way}, not just at arrival.}
\end{itemize}
Throughout this chapter we work in the fully observable,
reachability-goal case, where the differences between weak, strong, and
strong-cyclic plans are already visible.  richer combinations are
tackled by the same symbolic machinery, just with more expressive
logics --- exactly the logics we will meet in Part~II.
\end{formaldetails}

\begin{example}[Parcel delivery with a traffic light]
\label{ex:parcel-traffic}
We adopt the parcel-delivery domain used throughout the book.  A robot
must move a parcel from the train \emph{station} to the \emph{airport}.
Between them sits a traffic light, which may be \emph{green} or
\emph{red}.  The fluent $\mathit{pos}\in\{\mathit{station},\mathit{light},\mathit{airport}\}$
records where the robot is; the fluent $\mathit{color}\in\{\mathit{green},\mathit{red}\}$
records the light.  The actions are $\texttt{transport}$ (advance one
step along the route) and $\texttt{wait}$ (do nothing for one tick).
The position is inertial; the colour is non-inertial.  Crucially,
\texttt{transport} is allowed at the light only when the colour is
green, and a $\texttt{transport}$ from the station may land at the light
with either colour.  \cref{fig:parcel-domain} shows the resulting six
states.
\end{example}

\begin{figure}[ht]
\centering
\begin{tikzpicture}[
  >={Stealth[length=2.2mm]},
  st/.style={circle, draw, thick, minimum size=13mm, inner sep=0pt,
             font=\footnotesize, align=center},
  init/.style={st, fill=teal!10},
  plain/.style={st, fill=gray!8},
  goal/.style={st, draw, double, thick, fill=teal!16},
  nondet/.style={->, thick, draw=orange!75!black},
  det/.style={->, thick, draw=graphite!70},
]

  \node[init] (sg) at (0,1.3)  {$\mathsf{S}$\\$\mathit{green}$};
  \node[init] (sr) at (0,-1.3) {$\mathsf{S}$\\$\mathit{red}$};

  \node[plain] (lg) at (5,1.3)  {$\mathsf{L}$\\$\mathit{green}$};
  \node[plain] (lr) at (5,-1.3) {$\mathsf{L}$\\$\mathit{red}$};

  \node[goal] (ag) at (10,1.3)  {$\mathsf{A}$\\$\mathit{green}$};
  \node[goal] (ar) at (10,-1.3) {$\mathsf{A}$\\$\mathit{red}$};


  \draw[nondet, bend left=15]  (sg) to node[above, font=\scriptsize]{$\texttt{transport}$} (lg);
  \draw[nondet, bend right=15] (sg) to node[below, font=\scriptsize, yshift=-1mm]{$\texttt{transport}$} (lr);
  \draw[nondet, bend left=15]  (sr) to (lg);
  \draw[nondet, bend right=15] (sr) to (lr);



  \draw[det, bend left=18] (lr.north) to[bend left=18]
        node[right, font=\scriptsize, xshift=1mm]{$\texttt{wait}$} (lg.south);
  \draw[det, bend left=18] (lg.south) to[bend left=18]
        node[left, font=\scriptsize, xshift=-1mm]{$\texttt{wait}$} (lr.north);


  \draw[det] (lg) -- node[above, font=\scriptsize]{$\texttt{transport}$} (ag);
  \draw[det] (lg) to[bend left=22] (ar);


  \node[draw=graphite!30, dashed, rounded corners,
        fit=(ag) (ar), inner sep=3mm,
        label={[font=\scriptsize, color=graphite!70]above:goal $G$}] {};
  \node[draw=graphite!30, dashed, rounded corners,
        fit=(sg) (sr), inner sep=3mm,
        label={[font=\scriptsize, color=graphite!70]left:$I$}] {};
\end{tikzpicture}
\caption{The parcel-delivery domain.  Initial states $\mathsf{S}_\mathit{green},\mathsf{S}_\mathit{red}$
are shown in teal, goal states $\mathsf{A}_\mathit{green},\mathsf{A}_\mathit{red}$ are double-circled in teal.
The orange $\texttt{transport}$ edges from the station are non-deterministic: the environment
chooses the colour on arrival.  The grey $\texttt{wait}$ loop between $\mathsf{L}_\mathit{red}$ and
$\mathsf{L}_\mathit{green}$ is what will force us to consider strong-cyclic plans.}
\label{fig:parcel-domain}
\end{figure}

This little domain already contains every complication we need.  The
two initial states reflect uncertainty about the colour at the start.
The transition from the station is non-deterministic, because the
environment picks the colour on arrival.  And the wait-loop between
$\mathsf{L}_\mathit{red}$ and $\mathsf{L}_\mathit{green}$ means a naive
``try to cross, then wait, then try again'' policy can in principle
loop forever.  Each plan type in \cref{sec:plans-execution-structures}
will be defined so that it copes with one of these complications.


% ============================================================
\section{Plans as Execution Structures}
\label{sec:plans-execution-structures}
% ============================================================

A domain tells us what \emph{could} happen; a plan tells us what the
planner \emph{chooses} to do.  The two together generate an execution
structure, a subgraph of the domain in which only the chosen actions
survive but every environmental outcome of each chosen action is kept.

\begin{definition}[Plan as a state-action table]
\label{def:plan}
A \textbf{plan} (or \textbf{state-action table}) for a domain is a
relation $\pi\subseteq S\times A$ with $(s,a)\in\pi$ only when $a$ is
enabled at $s$.  The plan is \textbf{deterministic} if for every $s$
there is at most one $a$ with $(s,a)\in\pi$; it is \textbf{universal}
if for every $s\in S$ there is at least one such $a$.
\end{definition}

In a fully deterministic domain, a plan collapses to the familiar
object: a finite sequence $a_1,a_2,\ldots,a_k$ of actions, since at
each step there is only one enabled action and only one possible
successor.  In a non-deterministic domain this is no longer enough.
The fixed sequence may fail after an unlucky outcome, so the plan must
be allowed to branch and react --- exactly like a strategy in a game.

\begin{definition}[Execution structure]
\label{def:execution-structure}
Let $\pi$ be a plan and $I\subseteq S$ a set of initial states.  The
\textbf{execution structure} induced by $\pi$ from $I$ is the directed
graph $K_\pi=(Q,T_\pi)$ where $Q\subseteq S$ and $T_\pi\subseteq
Q\times Q$ are the minimal sets satisfying:
\begin{itemize}[nosep]
  \item $I\subseteq Q$;
  \item if $s\in Q$ and $(s,a)\in\pi$ and $(s,a,s')\in T$, then
        $s'\in Q$ and $(s,s')\in T_\pi$.
\end{itemize}
A node $s\in Q$ is \textbf{terminal} if it has no outgoing edge in
$T_\pi$.  An \textbf{execution path} from $s_0\in I$ is a finite or
infinite sequence $s_0,s_1,s_2,\ldots$ with $(s_i,s_{i+1})\in T_\pi$
for every $i$, ending in a terminal state if it is finite.  The
structure is \textbf{acyclic} if every execution path is finite.
\end{definition}

The definition makes a subtle but crucial move.  The plan $\pi$ lives
in the world of state-action pairs; the execution structure lives in
the world of state-state pairs.  Closing $\pi$ under all environmental
outcomes produces a graph that contains every future the planner cannot
rule out.  Terminal states are where this graph stops: either because
the planner reached its target, or because it reached a dead end with
no enabled action.

\begin{definition}[Weak, strong, and strong-cyclic solutions]
\label{def:weak-strong-plan}
\label{def:strong-cyclic-plan}
Let $K_\pi=(Q,T_\pi)$ be the execution structure induced by $\pi$ from
$I$, and let $G$ be the goal set.
\begin{itemize}[nosep]
  \item $\pi$ is a \textbf{weak solution} if, for every $s_0\in I$,
        there \emph{exists} an execution path from $s_0$ that ends in a
        terminal state belonging to $G$.
  \item $\pi$ is a \textbf{strong solution} if $K_\pi$ is acyclic and
        \emph{every} terminal state of $K_\pi$ belongs to $G$.
  \item $\pi$ is a \textbf{strong-cyclic solution} if from every state
        of $Q$ some terminal state of $K_\pi$ is reachable, and every
        terminal state of $K_\pi$ belongs to $G$.
\end{itemize}
\end{definition}

These three notions increase in strength.  Weak planning is
\emph{existential}: one lucky path suffices.  Strong planning is
\emph{universal and acyclic}: every path terminates, and every
termination is a goal.  Strong-cyclic planning sits between the two:
cycles are tolerated, but only if from every state along them a goal
remains reachable, and a fairness assumption rules out the run that
loops forever while postponing the exit.

\begin{example}[Weak, strong, and strong-cyclic on the parcel domain]
\label{ex:parcel-plan-types}
In \cref{fig:parcel-domain} no strong plan can exist.  Whatever the
planner does at $\mathsf{L}_\mathit{red}$, the environmental outcome
``stay red'' is always available, so any execution structure retains
the cycle $\mathsf{L}_\mathit{red}\to\mathsf{L}_\mathit{green}\to\mathsf{L}_\mathit{red}$.
A weak plan exists trivially: a path that crosses only when green
\emph{can} reach the airport in finitely many steps.  A strong-cyclic
plan also exists, because from every state in the structure an airport
state remains reachable, and under the fairness assumption that the
green outcome cannot be postponed forever, the robot eventually crosses.
This is exactly the situation safety-critical AI planning cares about:
we cannot \emph{force} success, but we can guarantee it under fair
scheduling.
\end{example}

\begin{definition}[Situated, universal, and conformant plans]
\label{def:planning-plan-types}
A \textbf{situated plan} assumes a known initial state.  A
\textbf{universal plan} gives an action choice for every reachable state
in a set of possible states.  A \textbf{conformant plan} must work
without observing which of several possible initial states is the real
one.
\end{definition}

\begin{remark}[Determinization]
\label{rem:determinization}
A non-deterministic plan $\pi$ is a weak (respectively strong,
strong-cyclic) solution if and only if \emph{every} deterministic
sub-plan $\pi'\subseteq\pi$ that still covers all states of $\pi$ is
also a weak (respectively strong, strong-cyclic) solution.  This
universal quantifier over determinizations is what lets us search the
smaller space of deterministic plans without loss of strength.
\end{remark}

Strong cyclic planning is useful precisely when recovery loops are
acceptable but strong guarantees are unachievable.  If the robot
repeatedly tries to cross while the light may nondeterministically be
red or green, a strong plan is impossible because the red outcome can
keep it stuck forever.  A strong cyclic plan is still reasonable under
the fairness assumption that the green outcome eventually occurs ---
the same kind of compassion constraint that we will revisit when we
study fair model checking in Part~II.

A useful mnemonic is the logic-flavour summary:

\begin{center}
\begin{tabular}{@{}lll@{}}
\toprule
Plan type & Path quantifier & Acyclicity \\
\midrule
weak        & $\exists$ path reaching $G$      & not required \\
strong      & $\forall$ paths, terminal in $G$ & required \\
strong-cyclic & $\forall$ states, $\exists$ path to $G$ (under fairness) & cycles allowed \\
\bottomrule
\end{tabular}
\end{center}


% ============================================================
\section{Preimage Operators}
\label{sec:preimage-operators}
% ============================================================

How do we compute, rather than just describe, the states from which the
planner can win?  The answer is the same trick that drove the attractor
construction in \cref{ch:synthesis-games}: work backwards from the goal,
one step at a time, and ask which states can be forced into the current
winning set by a single action choice.  The two flavours of plan ---
weak and strong --- give rise to two one-step predecessors, the
\emph{preimage operators}.

\begin{definition}[Weak preimage]
\label{def:weak-preimage}
The weak preimage of $X\subseteq S$ is
\[
  \mathsf{Pre}_w(X)=
  \{s\in S:\exists a\in A,\exists s'\in X
        \text{ such that }(s,a,s')\in T\}.
\]
\end{definition}

This is the one-step version of a weak plan: there exists an action and
an outcome that reaches $X$.  It captures the planner's ability to be
\emph{lucky} --- the environment is allowed to pick the friendliest
outcome of the chosen action.

\begin{definition}[Strong preimage]
\label{def:strong-preimage}
The strong preimage of $X\subseteq S$ is
\[
  \mathsf{Pre}_s(X)=
  \{s\in S:\exists a\in A,\;
        \mathrm{Post}(s,a)\neq\emptyset
        \text{ and }
        \mathrm{Post}(s,a)\subseteq X\}.
\]
\end{definition}

This is the one-step version of a strong plan.  The planner chooses the
action, but nature chooses the outcome.  Therefore every possible
successor must already lie inside $X$; one bad outcome is enough to
disqualify $s$ from being a strong predecessor.

\begin{figure}[ht]
\centering
\begin{tikzpicture}[
  >={Stealth[length=2.2mm]},
  state/.style={circle, draw, thick, minimum size=8mm, fill=gray!10},
  goal/.style={circle, draw, double, thick, minimum size=8mm, fill=teal!18},
  warr/.style={->, thick, draw=orange!80!black},
  sarr/.style={->, thick, draw=graphite!75},
]

  \node[state] (s) at (0,0.9)  {$s$};
  \node[goal]  (x1) at (3,1.7) {$x_1$};
  \node[goal]  (x2) at (3,0.1) {$x_2$};

  \node[state] (t) at (7,0.9)  {$t$};
  \node[goal]  (y1) at (10,1.7){$y_1$};
  \node[state] (y2) at (10,0.1){$y_2$};


  \draw[sarr] (s) -- node[above, font=\scriptsize]{$a$} (x1);
  \draw[sarr] (s) -- node[below, font=\scriptsize]{$a$} (x2);

  \draw[warr] (t) -- node[above, font=\scriptsize]{$a$} (y1);
  \draw[warr] (t) -- node[below, font=\scriptsize]{$a$} (y2);

  \node[below=4mm of x2, font=\scriptsize, align=center]
        {all outcomes in $X$\\strong predecessor};
  \node[below=4mm of y2, font=\scriptsize, align=center]
        {one outcome escapes\\only weak predecessor};
\end{tikzpicture}
\caption{Strong vs.\ weak preimage.  In graphite, $s\in\mathsf{Pre}_s(X)$:
all outcomes of action $a$ land in $X$.  In orange, $t\notin\mathsf{Pre}_s(X)$
but $t\in\mathsf{Pre}_w(X)$: one outcome escapes, so $a$ cannot \emph{force}
entry into $X$.}
\label{fig:strong-preimage}
\end{figure}

\Cref{fig:strong-preimage} makes the operational difference explicit.
The state $s$ is a strong predecessor of $X$ because action $a$ pins
both successors inside $X$.  The state $t$ is not: even though $a$ has
one successor in $X$, it has another outside, and an adversarial
environment can pick that one.  This is precisely the same
``controllable predecessor'' operator that built the system attractor
in the games of \cref{ch:synthesis-games}; only the names have changed.


% ============================================================
\section{Backward Fixpoint Algorithms}
\label{sec:planning-fixpoints}
% ============================================================

With the preimage operators in hand, the strong and weak planning
problems reduce to least-fixpoint computations on the lattice
$(2^S,\subseteq)$.  The goal set is the seed; we repeatedly add states
from which the planner can force progress into the current winning set;
when nothing new appears, we have reached the fixpoint.

\begin{formaldetails}[Strong backward planning]
\begin{enumerate}[label=\textbf{Step \arabic*.}, leftmargin=2cm]
  \item Initialize $W_0=G$.
  \item Compute
        \[
          W_{i+1}=W_i\cup \mathsf{Pre}_s(W_i).
        \]
  \item Stop when $W_{i+1}=W_i$.
  \item The resulting fixpoint $W^\ast$ is the set of states from which
        the goal can be \emph{forced}.
\end{enumerate}
The weak variant is identical except that $\mathsf{Pre}_s$ is replaced
by $\mathsf{Pre}_w$.
\end{formaldetails}

This is the same logical shape as the attractor computation in
\cref{ch:synthesis-games}.  At each step we add states from which the
controller has an action whose possible outcomes are already known to
be winning.  The iteration terminates because $S$ is finite and the
sequence is monotone.  By the Knaster--Tarski theorem the result is
exactly the least fixpoint of the operator $X\mapsto G\cup
\mathsf{Pre}_s(X)$.

\begin{proposition}[Strong winning region]
\label{prop:strong-planning-fixpoint}
Let $W^\ast$ be the fixpoint above.  If an initial state belongs to
$W^\ast$, then a strong plan exists from that state.  If it does not
belong to $W^\ast$, no memoryless strong plan based on the represented
state space can force the goal using the available transitions.
\end{proposition}

\begin{proof}
For the first claim, assign a rank to each state according to the first
iteration in which it enters the sequence $W_0,W_1,\ldots$.  States of
rank $0$ are already goals.  If a state has rank $i+1$, then by
definition it belongs to $\mathsf{Pre}_s(W_i)$, so there exists an
action whose every successor lies in $W_i$.  Choose such an action.  No
matter which successor occurs, the rank strictly decreases.  Since ranks
are natural numbers, every execution reaches rank $0$, hence a goal.

For the second claim, let $W^\ast$ be the fixpoint and consider a state
$s\notin W^\ast$.  For every available action, either it has no valid
successor or it has some successor outside $W^\ast$; otherwise $s$
would belong to $\mathsf{Pre}_s(W^\ast)$ and hence to $W^\ast$ by the
fixpoint property.  An adversarial environment can always choose such
an outside successor, keeping the execution outside $W^\ast$ forever.
Therefore no state-based strong plan can force the goal from outside
the fixpoint.
\end{proof}

The extracted policy chooses, at each newly added state, one action
that witnessed membership in $\mathsf{Pre}_s(W_i)$.  Ranks double as a
certificate: they bound the number of steps to the goal.

\begin{formaldetails}[Strong-cyclic backward planning]
Strong-cyclic planning sits between weak and strong planning.  A weak
plan only needs one successful outcome; a strong plan must tolerate all
outcomes.  A strong-cyclic plan tolerates all outcomes \emph{provided}
the chosen policy never leaves a closed live region and fairness rules
out postponing enabled progress forever.

The standard computation therefore alternates two requirements.  Given
a candidate region $R$, keep a non-goal state $s\in R$ only if there is
an action $a$ such that
\[
  \mathsf{Post}(s,a)\neq\emptyset,\qquad
  \mathsf{Post}(s,a)\subseteq R,\qquad
  \mathsf{Post}(s,a)\cap X\neq\emptyset
\]
for the current progress set $X$.  The first two clauses say that the
policy is enabled and closed inside the candidate region; the third says
that, under a fair resolution of nondeterminism, the policy has a way to
make progress toward states already known to be closer to the goal.
Starting from $X=G$ and iterating this progress test inside a greatest
closed candidate yields the usual nested-fixpoint view of
strong-cyclic planning.

Operationally, one can read the algorithm as a backward pass followed by
pruning, but the pruning is not optional bookkeeping: it removes states
for which every apparent weak path to the goal depends on first leaving
the live region or on cycling forever without any fair progress edge.
The resulting policy has two invariants: every possible successor stays
inside the region, and from every reachable non-goal state the goal is
again reachable through policy edges.
\end{formaldetails}

This is the precise content behind the slogan ``from every state, a
goal remains reachable, and every terminal state is a goal''.  The
policy is not allowed to gamble on a single lucky branch: it must keep
all nondeterministic outcomes inside the controlled region.  Fairness
then supplies the missing ingredient that distinguishes strong-cyclic
planning from strong planning: if progress remains continuously
available, an execution cannot ignore it forever.


% ============================================================
\section{Symbolic Representation}
\label{sec:symbolic-planning}
% ============================================================

Explicit state graphs can be enormous.  The parcel domain has six
states, but a domain with forty Boolean fluents already hides up to
$2^{40}$ states, and most realistic planning problems have many more.
Symbolic model checking sidesteps the explosion by representing sets
and transitions with Boolean formulas and manipulating whole sets at
once.  The preimage operators we defined set-theoretically translate
cleanly into quantified Boolean formulas, where they can be evaluated by
BDD-based or SAT-based engines.

\begin{definition}[Symbolic transition relation]
\label{def:symbolic-transition}
Let $x$ be the vector of current-state variables, $a$ the action
variables, and $x'$ the next-state variables.  A symbolic transition
relation is a Boolean formula
\[
  T(x,a,x')
\]
that is true exactly for valid transitions.
\end{definition}

Why represent the chosen action with a vector of Boolean variables $a$
rather than a single scalar?  As the professor noted, assigning an
independent Boolean variable to each action naturally models
\textbf{concurrent actions}: multiple action variables can be true at
once.  If the domain forbids concurrent execution, the modeller must
explicitly add a mutual-exclusion constraint (e.g., $\sum a_i \le 1$) to
the transition relation.  When concurrency is ruled out structurally, a
more compact scalar encoding (using $\lceil\log_2|A|\rceil$ bits)
suffices.

A set of states is represented by a formula $X(x)$.  The strong preimage
can then be written logically as
\[
  \mathsf{Pre}_s(X)(x)
  =
  \exists a\;
  \Big(
    \exists x'\,T(x,a,x')
    \;\wedge\;
    \forall x'\,(T(x,a,x')\rightarrow X(x'))
  \Big).
\]
Read this formula carefully: the outermost $\exists a$ is the planner's
choice of action, the inner $\exists x'\,T$ checks that the action is
actually enabled, and the universal $\forall x'$ is nature's adversarial
choice of outcome, which we force to land inside $X$.

The weak preimage has the existential version of the same shape:
\[
  \mathsf{Pre}_w(X)(x)
  =
  \exists a\,\exists x'\,
  \big(T(x,a,x')\wedge X(x')\big).
\]
The strong preimage replaces the existential choice of outcome with a
universal requirement over all outcomes of the chosen action.  In
implementation, these formulas are rarely evaluated by enumerating
assignments: they are compiled into ordered binary decision diagrams
(OBDDs), on which projection and Boolean operations are polynomial in
the diagram size.  Part~II develops this symbolic machinery in detail.

\begin{intermezzo}{Why Symbolic Methods Matter}
If a domain has 40 Boolean fluents, the explicit state space may contain
$2^{40}$ states.  Symbolic methods do not make the problem magically
small, but they let us manipulate large regular sets of states at once,
often using compact representations such as ordered binary decision
diagrams.  Many planning domains are mostly regular: a robot's dynamics
factor over its fluents, and a single BDD node can summarise millions
of concrete states that all behave identically.
\end{intermezzo}


% ============================================================
\section{Action Description Languages}
\label{sec:action-description-languages}
% ============================================================

Writing the transition relation $T(x,a,x')$ by hand is tedious and
error-prone.  Planning tools therefore let the modeller declare, at a
higher level, what each action \emph{requires} and what it
\emph{causes}; a compiler then expands those declarations into the
symbolic transition relation used by the fixpoint algorithms.  This is
the planning analogue of writing a program in a high-level language and
letting a compiler lower it to machine code.

\begin{definition}[Action rule]
\label{def:action-rule}
An action rule describes when an action is executable and how it changes
fluents.  It typically has a precondition part and an effect part.
\end{definition}

A typical rule has the form
\[
  A \text{ causes } P \text{ if } Q,
\]
meaning that when action $A$ is executed in a state satisfying $Q$, the
next state must satisfy effect $P$.  Executability conditions say when
an action is allowed; frame or inertia constraints say that inertial
fluents keep their current value unless some rule changes them.

\begin{example}[Compiling one action rule]
\label{ex:compile-action-rule}
The rule
\[
  \texttt{load}\text{ causes }\mathit{loaded}
  \text{ if }\mathit{robotAtParcel}
\]
contributes the implication
\[
  \texttt{load}(a)\wedge \mathit{robotAtParcel}(x)
  \;\Longrightarrow\;
  \mathit{loaded}(x')
\]
to the transition relation.  Inertia contributes, for example,
\[
  \neg\textsf{changesLoaded}(x,a)
  \;\Longrightarrow\;
  \big(\mathit{loaded}(x')\leftrightarrow \mathit{loaded}(x)\big).
\]
\end{example}

\begin{example}[Rules for the parcel domain]
\label{ex:parcel-rules}
The parcel domain of \cref{fig:parcel-domain} can be written
declaratively as follows.  State invariants pin down what a state is:
\[
  \textsf{always}\;\Big(\mathit{pos}=\mathit{station}\oplus
  \mathit{pos}=\mathit{light}\oplus\mathit{pos}=\mathit{airport}\Big).
\]
Initial and goal states are declared with the predicates
\textsf{initially} and \textsf{goal}:
\[
  \textsf{initially}\;\big(\mathit{pos}=\mathit{station}\big),
  \qquad
  \textsf{goal}\;\big(\mathit{pos}=\mathit{airport}\big).
\]
The actions are compiled into transition constraints:
\[
  \texttt{transport}\;\textsf{causes}\;
  \mathit{pos}=\mathit{light}\;\textsf{if}\;
  \mathit{pos}=\mathit{station},
\]
\[
  \texttt{transport}\;\textsf{causes}\;
  \mathit{pos}=\mathit{airport}\;\textsf{if}\;
  \mathit{pos}=\mathit{light}\wedge \mathit{color}=\mathit{green}.
\]
The colour $\mathit{color}$ is non-inertial, so its evolution is left
unconstrained by these rules; nature chooses it.  A small set of
declarative rules has thus generated the explicit six-state graph of
\cref{fig:parcel-domain}.
\end{example}

% ============================================================
\section{Towards Model Checking}
\label{sec:planning-towards-mc}
% ============================================================

The chapter has traced a single arc.  Planning problems are finite
transition systems with goals.  Plans are state-action tables, and
their strength is measured against the universal or existential
path-quantification in the induced execution structure.  Strong and
weak solutions are characterised by preimage fixpoints; strong-cyclic
solutions add a forward reachability pass on top.  The whole pipeline
compiles declaratively from an action description language into the
symbolic transition relations of model checking:

\[
\begin{tikzpicture}[baseline=-0.6ex,
  >={Stealth[length=2mm]},
  node distance=14mm and 9mm,
  box/.style={draw, rounded corners=2pt, fill=teal!8,
              minimum height=8mm, inner sep=4pt, font=\small, align=center}]
  \node[box] (a) {action\\rules};
  \node[box, right=of a] (b) {symbolic\\transition $T$};
  \node[box, right=of b] (c) {preimage\\fixpoint};
  \node[box, right=of c] (d) {weak / strong /\\strong-cyclic\\winning region};
  \draw[->, thick] (a) -- (b);
  \draw[->, thick] (b) -- (c);
  \draw[->, thick] (c) -- (d);
\end{tikzpicture}
\]

The closing observation is that the weak and strong plan-existence
questions are themselves expressible in branching-time temporal logic.
The existence of a path to a goal is the CTL formula
$\mathit{EF}\,G$; the universality of reaching a goal is
$\mathit{AF}\,G$; the strong-cyclic guarantee, under fairness, is a
fair variant of $\mathit{AF}\,G$.  In other words, the very same CTL
that will drive the symbolic model checkers of Part~II already
classifies the plans of this chapter.  Planning \emph{is} model
checking on a finite transition system, with a reachability objective.

\begin{summarybox}
We have now closed Part~I.  Starting from finite words and regular
languages, we added infinite words, complementation, infinite games,
and finally planning.  The throughline has been the same: a finite
transition system, a target set, and a preimage-based backward search.
\begin{itemize}[nosep]
  \item Planning domains are finite transition systems with actions.
  \item Plans are state-action tables; under non-determinism they
        become strategies.
  \item Weak planning is existential; strong planning is universal;
        strong-cyclic planning adds fairness on top of reachability.
  \item Preimage operators are one-step controllable predecessors,
        reused from the attractor construction of \cref{ch:synthesis-games}.
  \item Symbolic model checking represents large planning problems
        with Boolean formulas and computes their fixpoints without
        enumerating states.
\end{itemize}

Part~I has built the theoretical machinery: finite and infinite
automata, games, and the attractor/preimage reasoning that solves them.
Part~II turns to scale and practice.  We will introduce temporal logics
(LTL and CTL) as specification languages, study explicit and symbolic
model-checking algorithms, and meet the toolchain (SMV, OBDDs,
SAT-based bounded model checking, learning, and Tamarin) that turns
the theory of these chapters into running verifiers.  Many of the
questions Part~I leaves open --- how do we specify temporal goals
beyond reachability, how do we handle fairness explicitly, how do we
verify systems far larger than six states --- are exactly what Part~II
is designed to answer.
\end{summarybox}

\section*{Exercises}
\addcontentsline{toc}{section}{Exercises}

\begin{chapex}[Weak vs. strong on the parcel domain]
In \cref{fig:parcel-domain}, list the contents of $\mathsf{Pre}_w(G)$
and $\mathsf{Pre}_s(G)$ for $G=\{\mathsf{A}_\mathit{green},\mathsf{A}_\mathit{red}\}$
after a single iteration.  Explain in one sentence why
$\mathsf{L}_\mathit{red}$ belongs to the first set but not the second.
\end{chapex}

\begin{chapex}[One fixpoint step]
Let $G=\{g\}$ and suppose state $s$ has an action whose only successors
are $g$ and another state already in $G$.  Does $s$ enter the strong
winning set after one iteration?  Justify by evaluating the membership
condition for $\mathsf{Pre}_s(G)$.
\end{chapex}

\begin{chapex}[Strong, weak, or strong-cyclic?]
Consider the parcel domain of \cref{fig:parcel-domain} and the plan
``at the light, if green then \texttt{transport}, else \texttt{wait};
at the station, \texttt{transport}''.  Classify it as weak, strong, or
strong-cyclic, and explain which feature of \cref{fig:parcel-domain}
rules out the stronger classifications.
\end{chapex}

\begin{chapex}[Symbolic preimage]
In the symbolic strong preimage formula
\[
  \exists a\;\big(\exists x'\,T(x,a,x')\wedge \forall x'\,(T(x,a,x')\rightarrow X(x'))\big),
\]
identify which quantifier is controlled by the planner, which expresses
action executability, and which expresses adversarial outcome
selection by the environment.
\end{chapex}
